\section{RESEARCH PLAN}
\subsection{Motivation}

Traditionally, watching a television program has been a passive experience:
All viewers receive the same TV content at the same
playback pace, regardless of individual preferences.
Not being able to select viewing pace 
can be inconvenient and even unproductive 
sometimes. For instance, when one
is interrupted while watching a news program
(say, by a phone call or biological need), one may 
miss the information (e.g., stock quotes, weather forecast,
etc.) that one has been waiting for.  Similarly, when
one misses a crucial point while viewing a televised
lecture, one may not be able to comprehend the materials
after that point.
What is desirable is {\em viewing interactivity}
---pause, fast-forward and instant replay---allowing
the time-shifting of a program so the viewer can
watch at whatever pace he or she finds comfortable.
One can also record programs at multiple channels at the same
time (which a tape-based VCR cannot do), filter out unwanted
content and produce a customized program. 

In addition to viewing interactivity, being able to 
interact with the content of a program (what I call
{\em content interactivity}) can enrich viewing experience.
Content interactivity is
especially beneficial for educational programs. For instance,
children would be able to change the attributes 
(e.g., colors, size, position 
and even the personality) of the ``bunny,'' ``big bird'' and
other characters in the {\it Sesame Street} program.
A grade school student can learn American states
by interacting with a map.
A class of high school students could conduct science experiments
in small groups through a virtual lab and their instructor
could keep track of their progress in real-time. 
Many other applications exist.

To support interactivity, many studies suggest a server-based
architecture (e.g., \cite{ex80,ex49,ex54,ex1,ex86}),
in which a server schedules
its resources (i.e., memory and disk bandwidth)
to support interactivity.
In a server-based architecture, the receivers are assumed
to be passive---simply receiving bits and rendering frames.
However, because of the typical long end-to-end transmission
delay between a server and a receiver and the Interet's limited bandwidth,
it is practically infeasible for the server
to support ``real-time'' viewing interactivity for tens of
thousands of simultaneous users.  Furthermore,
on a broadcast channel (e.g., CNN), one simply cannot request
the server to pause or replay a program.
 
Instead of the server-based approach, I plan to investigate
new hybrids of technologies for 
a receiver-based architecture to support viewing interactivity 
and content interactivity. Figure~\ref{fig-dtv-pipeline} depicts a 
multimedia data delivery pipeline.  A receiver in this architecture
is equipped with a large and inexpensive disk and may have an 
Internet back-channel in addition to the broadcast channel.
The economical caching together with the random
access capability of the disk enables a multimedia receiver to 
buffer a large amount of data to support
time-shift operations.  A viewer can pause a live 
program to take a break from viewing while the broadcast
stream continues arriving and being written to the local disk.
The viewer can resume watching the program after the pause
with a delay, or fast-forward the program to get back in synch with
the broadcast stream.  
While a given broadcast stream remains the same
for all viewers, various data objects can be transmitted to
individual viewers via the Internal channel.
Image-based rendering techniques can
overlay a number of data objects with the broadcast stream to
produce a customized program for content
interaction.

\begin{figure}[t]
\epsfxsize=350pt
\centerline{\epsffile{fig-dtv-pipeline.eps}}
\caption{The Multimedia Data Delivery Pipeline.}
\label{fig-dtv-pipeline}
\vspace{.1in}
\end{figure}

Extending the DRAM cache with a large disk and
adding an Internet channel to the receiver add complexity to 
the receiver's buffer and data management. To ensure that IO 
are efficiently scheduled, the receiver needs a 
{\em buffer manager} that can schedule the
receiver's memory and
disk bandwidth effectively to service
a wide variety of interactive scenarios.
The receiver also needs a {\em data manager}
that can manage data from TV, radio and the Web
for supporting efficient queries and data
retrieval. Although progress has been made in these areas,
many issues, especially in managing multimedia
data and managing buffer for rendering and
for querying multimedia data,
and in developing an architecture that can be adaptive
to both resource constraints and {\em Quality of
Service} (QoS) requirements, remain open.
(These challenges are discussed
in detail in Section~\ref{sec-challenges}.)
I plan to address these issues, and to build on
work from several fields, most notably from the
image processing, operating system and database fields.
In particular, I am interested in developing the
following technologies:
                             
\begin{nlist}
\item {\bf 1}.
Temporal indexing structures that support
efficient stream synchronization.
\item {\bf 2}. Caching and prefetching techniques that satisfy
the real-time data requirement of the continuous media data
(i.e., video and audio), on the one hand, and minimize the
initial latency in responding to an interactive request,
on the other.
\item {\bf 3}.
QoS-adaptive buffer-management architectures that manage
memory and disk bandwidth adaptively to
the receiver's resource constraints and
the viewer's QoS requirements.
\item {\bf 4}.
Media query processing that handles heterogeneous
media queries efficiently.
\end{nlist}   

In addition, I envision constructing demonstrations
that use these technologies.
Here is a list of some of the demonstrations I am planning
to implement:
\begin{jlist}
\item 
A multimedia database that supports content-based
queries,
\item 
A digital VCR that supports pause, fast-forward and instant
replay, and
\item 
A multi-camera video surveillance system (e.g., for
monitoring in a nuclear power facility).
\end{jlist}

The multimedia database system provides the digital-VCR
project with an efficient storage and an effective query
engine that can help the viewer select content
before the viewer issues a fast-forward or replay command.
The digital VCR, which receives only one stream at a time,
in turn, is extended by the multi-camera video surveillance project,
which handles multiple cameras, streams and viewers simultaneously.
These demonstrations show that 
a wide variety of interactive multimedia applications
in education, commercial, government and
military sectors can be built on the technologies
developed in this project.
                           
Much of the research I propose will be joint efforts
with colleagues and students at the University of California,
Santa Barbara and Stanford University.

In the remaining sections of this proposal, I enumerate the challenges
this project faces, the core technologies I plan to develop, 
my prior accomplishments, and
the demonstrations I plan to implement.   

\subsection{Technical Challenges}
\label{sec-challenges}

A multimedia receiver faces at least three design challenges:
\begin{jlist}
\item Not all data are equal: different types of data may have different
size and latency characteristics.
\item Not all receivers are equal: different receivers may have different
available hardware resources and network bandwidth.
\item Not all viewers are equal: different viewers may have different
viewing preferences.
\end{jlist}
In the remainder of this section, I enumerate these challenges.

\subsubsection{Not All Data Are Equal}

\begin{figure}[t]
\epsfxsize=260pt
\epsfysize=160pt
\centerline{\epsffile{fig-dtv-data.eps}}
\caption{The Characteristics of Multimedia Data Objects.}
\label{fig-dtv-data}
\end{figure}

Multimedia data objects can be continuous, textual, binary,
and so forth.  These data objects are
best characterized by their {\it sizes}
and {\it latency constraints}. Continuous       
data are audio/video streams, which are voluminous and
delay sensitive.
Textual data include captions, HTML and XML documents,
etc. Textual data occupy far less space than continuous data
and are not as delay sensitive. Other data, such as images and applets,
can be very copious or scant and may or may not be delay sensitive,
depending on the nature of the application.
Figure~\ref{fig-dtv-data} shows the characteristics of
these data objects. The $x$-axis in the figure represents the
delay sensitivity from low (on the left-hand side) to high
(on the right-hand side); the $y$-axis gives the size of an object.
At one extreme, videos are large and very sensitive to
delay. At the other extreme, captions are small in size and not very
sensitive to delay.
                       
\begin{figure}[t]
\epsfysize=160pt
\epsfxsize=280pt
\centerline{\epsffile{fig-delay-fs.eps}}
\caption{Delay Functions.}
\label{fig-delay-fs}
\end{figure}

When allocating memory and disk bandwidth for data objects, both 
their sizes and latency constraints must be taken into consideration.
We plan to characterize each data object's size and latency 
characteristic as follows:
The size of a data object is measured by the 
amount of storage space it needs.
The delay sensitivity of a data object is quantified as
a value function.  
As Figure~\ref{fig-delay-fs} shows, 
the function can be defined in many ways
(the $x$-axis in the figure depicts delay; the $y$-axis
depicts the value of the object normalized to
from zero i.e., worthless, to one): Function $F_a$
depicts an object that can tolerate delay up to a period $t_\delta$.
Function $F_b$ shows the value of an object that decays linearly
after $t_\delta$. Functions $F_c$ and $F_d$ have other
value decay patterns.                                 
In terms of data streams,
function $F_a$ is a good value function for an audio frame,
function $F_b$ may be a good value function for a video
frame since slight delay of a frame display may be tolerable,
function $F_c$ may be a good value function for a subtitle, and
function $F_d$ may be a good value function for a Web page.
                                                       
\subsubsection{Not All Receivers Are Equal}

The critical resources of a receiver include memory,
network bandwidth and disk bandwidth.
These resources can also be characteristized by their
capacity and latency. The memory
size and the disk size
are both constants that
can be obtained from the system configuration file.
The latency of a receiver
is best characteristized by the IO time between memory
and the local disk since this latency dominates the latency
in the storage hierarchy. The IO time for reading/writing
data object $i$ (denoted as $\tau_i$) can be written as:
$\tau_i = \gamma(d) + S_i/TR$,
where $\gamma(d)$ is the seek time and rotational
delay of the disk given $d$ tracks to seek before transferring
data, $S_i$ is the size of the $i^{th}$
data object, and $TR$ is the transfer rate of the disk.
Note that the latency just described
assumes that no IO conflict exists. In reality, the          
more objects are accessed, the longer the delay that
can be experienced by an IO.  
Scheduling conflicts must be taken
into consideration and must be resolved
in an order that achieves maximum total QoS
for the viewer.
                                             
\subsubsection{Not All Viewers Are Equal}

To measure a viewer's satisfaction, we must allow the
user to have a say about what is important and
what is not. We thus associate each data object
with a QoS parameter. A viewer can
assign a QoS value to each data object, either
implicitly or explicitly.
For instance, if a user decides to turn off
all interactive features, the resource manager can
assign a QoS factor of zero to data objects
that are needed to support interactive features.
Regarding the data objects that are
needed to support the playback, higher QoS can be assigned to
mission-critical data objects, such as broadcast
video and audio streams and lower QoS can be assigned
to optional data objects such as captions and applets.

